
% ============================
% CHAPTERS 21--25 FOR THE TOME
% Phase Thermodynamics and Sub-Planckian Structure
% ============================

\part{Phase Thermodynamics and Sub-Planckian Structure}

\chapter{Sub-Planckian Phase Order and Gradient Dynamics}

\section{Phase Order Below the Planck Scale}

In MK--43 Ultra, the fundamental phase field $\tau(x)$ admits sub-Planckian structure described by gradients
\[
\nabla_\mu \tau,\qquad \nabla_\mu \nabla_\nu \tau,
\]
whose fluctuations encode local time-shell deformation.
These gradients are well-defined even in regimes where classical spacetime is not, enabling the model to describe ordering far beneath $\ell_P$.

\section{Gradient Dynamics and Stability}

Define the phase-gradient magnitude
\[
\Phi^2 = g^{\mu\nu} \nabla_\mu \tau\, \nabla_\nu \tau.
\]
Sub-Planckian ordering stabilises when the local cost functional
\[
\mathcal{L}_\tau = \alpha\, \Phi^2 + \beta\, (\Box \tau)^2 + \gamma\, (\nabla \tau)^4
\]
is minimised.
The coefficients $\alpha,\beta,\gamma$ are emergent from ER-foam couplings, ensuring that phase order inherits microstructure from entanglement patterns.

\chapter{Entropy Currents on the Time-Phase Shell}

\section{Entropy as a Phase-Gradient Flow}

Define a local entropy density $s(x)$ associated with coarse-grained ER links and chromatic fluctuations.
Entropy current is defined by
\[
J^\mu_s = s\, u^\mu + \kappa\, \nabla^\mu \tau,
\]
combining bulk motion $u^\mu$ with intrinsic phase-gradient drift.
In MK--43 Ultra, entropy is transported not only by matter/energy flux but by the geometry of the time-shell itself.

\section{Conservation and Production}

The entropy equation becomes
\[
\nabla_\mu J^\mu_s = \sigma \ge 0,
\]
where $\sigma$ arises from ER reconnections and chromatic diffusion.
This term encodes \emph{geometric} entropy production at the sub-Planck scale.

\chapter{Thermal Signatures of ER-Foam and Chromatic Modes}

\section{Temperature from ER Connectivity}

Thermal fluctuations are characterised by ER-link variance:
\[
T \propto \langle (W_{ij} - \langle W\rangle)^2 \rangle.
\]
Regions with highly fluctuating ER connectivity exhibit higher effective temperature, linking microscale network disorder to macroscopic thermodynamic observables.

\section{Chromatic Contributions}

Each chromatic band contributes a thermal signature
\[
T_c = \lambda_c \langle (\chi_c - \langle\chi_c\rangle)^2\rangle,
\]
where $\lambda_c$ depends on the white-state proximity.
Near $\Wcore$, chromatic variance collapses, raising effective temperature and driving thermalisation.

\chapter{Phase-Driven Irreversibility and the Arrow of Time}

\section{Phase Monotonicity}

Time orientation emerges from monotonic increase of integrated phase:
\[
\mathcal{T}(x) = \int \Phi\, d\Sigma,
\]
which cannot decrease due to ER-foam entropy production.
Thus the arrow of time emerges from geometric and informational irreversibility.

\section{Irreversibility from ER Reconnection}

ER foams constantly restructure.
Reconnection events are not microscopically reversible due to information loss into chromatic mixing.
This produces a preferred temporal orientation even at the sub-Planck level.

\chapter{Equilibrium, Fluctuations, and Phase-Stability Bands}

\section{Phase-Stability Bands}

Define stability bands by inequalities on the curvature of $\tau$:
\[
\Delta_i: \qquad a_i \le \Box \tau \le b_i.
\]
Within each band, fluctuations settle into metastable configurations defining local "phases" of the time-shell.

\section{Phase Transitions}

Crossing bands triggers phase transitions driven by chromatic rebalancing or ER density redistribution.
These transitions define:
\begin{enumerate}
\item thermal bursts,
\item chromatic realignments,
\item local white-state precursor formation.
\end{enumerate}

\section{Equilibrium Conditions}

Local equilibrium occurs when
\[
\nabla_\mu J^\mu_s = 0,\qquad \nabla_\mu \tau = \mathrm{const},
\]
signalling vanishing entropy production and stabilised time-shell curvature.

